% titel: Blatt 8, A8: MgO-Elementarzelle
% Teil II: der durchgerechnete Weg.

\subsection*{Gegeben und gesucht}
Elementarzelle von MgO (Abbildung), Kantenlänge
$a = 4.21\cdot10^{-10}\,$m. Gesucht: Strukturtyp, Stöchiometrie-Nachweis,
Koordinationszahlen, Dichte, Vergleich der Gitterenergien MgO/NaF.

\subsection*{a) Strukturtyp}
Steinsalz- bzw.\ NaCl-Typ (kubisch-flächenzentriertes Anionengitter, Mg$^{2+}$
in den Oktaederlücken).

\subsection*{b) Stöchiometrie}
Je Zelle:
\begin{itemize}
  \item O$^{2-}$ auf den Flächen: $6\cdot\frac{1}{2} = 3$, in den Ecken:
        $8\cdot\frac{1}{8} = 1$ $\Rightarrow$ insgesamt 4.
  \item Mg$^{2+}$ in den 12 Kanten: $12\cdot\frac{1}{4}=3$, in der Mitte: 1
        $\Rightarrow$ insgesamt 4.
\end{itemize}
Verhältnis 4$:$4 $= 1$:1, also MgO erfüllt.

\subsection*{c) Koordinationszahlen}
Im Steinsalz-Gitter ist jedes Ion oktaedrisch von 6 Gegenionen umgeben.
$\Rightarrow$ KZ = 6 für Ion in der Wurffmitte, für das Flächenion **und**
für das Ecken-Ion.

\subsection*{d) Dichte}
Zeile $Z=4$.
\[
  \rho = \frac{Z\,M}{N_A\,a^3}
       = \frac{4\cdot 40.30}{6.022\cdot10^{23}\cdot(4.21\cdot10^{-8}\,\text{cm})^3}
       \approx 3.59\ g\,cm^{-3}.
\]

\begin{ergebnis}
a) NaCl-/Steinsalz-Typ. b) 4 MgO pro Zelle, Verhältnis 1:1 erfüllt.
c) KZ = 6 (Mitte, Fläche, Ecke). d) $\rho \approx 3.6\ g\,cm^{-3}$.
\end{ergebnis}

\subsection*{e) Gitterenergie MgO vs.\ NaF}
MgO ($-3850\ kJ\,mol^{-1}$) liegt deutlich unter NaF ($-929\ kJ\,mol^{-1}$):
\begin{itemize}
  \item \textbf{Ladungen:} Mg$^{2+}$/O$^{2-}$ (doppelt geladen) gegenüber
        Na$^+$/F$^-$ (einfach) $\Rightarrow$ Coulomb-Anziehung $\sim
        q_1q_2$ wird um den Faktor 4 grösser.
  \item \textbf{Radien:} Mg$^{2+}$ und O$^{2-}$ sind klein, der Abstand
        zwischen den Ionen ist kurz, was die Anziehung zusätzlich stärkt.
\end{itemize}
Damit ist die Gitterenergie von MgO viel negativer (stabileres Gitter).

\begin{kontrolle}
Die Dichte (3.6 g/cm$^3$) liegt im typischen Bereich ionischer Oxide.
Gitterenergien: doppelt geladene Ionen und wenig Ladungsabstand geben den
erwarteten starken Sprung zwischen MgO und NaF.
\end{kontrolle}

\begin{merksatz}
Elementarzelle: Z über Wurffaktoren zählen, Dichte mit
$\rho=ZM/(N_A a^3)$. Gitterenergie steigt mit Ladungsprodukt und sinkender
Ionengrösse.
\end{merksatz}
